\chapter{Master Problem Bank: Concentration Cells \& Liquid Junction Potential}
\label{chap:qb_concentration_cells}

\begin{tcolorbox}[enhanced,colback=subtleblue,colframe=primarynavy,arc=2mm,boxrule=1pt,
    title=\textbf{\large Part IV Organization: Concentration Cells \& Liquid Junction Potential}]
This part compiles all questions and quantitative problems on electrode concentration cells (gas \& amalgam), electrolyte concentration cells without transference, concentration cells with transference, and the liquid junction potential ($E_{lj}$) from all 8 books. Identical and redundant variants from other books are co-located beneath each primary problem with full source citations.
\end{tcolorbox}

\section{Subtopic 4.1: Electrode Concentration Cells (Gas \& Amalgam)}

\begin{problembox}
\textbf{Q.4.1} \hfill \textbf{[Primary Source: Neeraj Kumar Part 1 Ex-I Q95 / GRB  12.24]}
\label{q:qb-c4-01}

Calculate the potential at $25^\circ\text{C}$ of the gas concentration cell:
\begin{equation*}
    \ce{Pt, H2}(P_1 = 10\,\text{atm}) \mid \ce{HCl}(0.1\,\text{M}) \mid \ce{H2}(P_2 = 1\,\text{atm})\ce{, Pt}
\end{equation*}
\begin{multicols}{4}
\begin{enumerate}[label=(\alph*)]
    \item $+0.0296\,\text{V}$
    \item $-0.0296\,\text{V}$
    \item $+0.0591\,\text{V}$
    \item $0.000\,\text{V}$
\end{enumerate}
\end{multicols}

\tcbline
\textbf{\color{accentamber}Cross-Book Redundant / Identical Variants:}
\begin{itemize}[leftmargin=*,itemsep=2pt]
    \item \textbf{[Variant v1: Pearson Ex-I Q90, p. 8]} For the gas cell $\ce{Pt, Cl2}(P_1) \mid \ce{NaCl}(1\,\text{M}) \mid \ce{Cl2}(P_2)\ce{, Pt}$, the cell reaction is spontaneous if: (a) $P_1 > P_2$ (b) $P_2 > P_1$ (c) $P_1 = P_2$ (d) independent of pressure.
    \item \textbf{[Variant v2: Narendra Avasthi Level-2 Q55, p. 30]} For a zinc amalgam concentration cell $\ce{Zn(Hg, } a_1\ce{)} \mid \ce{ZnSO4(aq)} \mid \ce{Zn(Hg, } a_2\ce{)}$, the EMF is given by: (a) $\frac{RT}{2F}\ln \frac{a_1}{a_2}$ (b) $\frac{RT}{2F}\ln \frac{a_2}{a_1}$ (c) $\frac{RT}{F}\ln \frac{a_1}{a_2}$ (d) zero.
    \item \textbf{[Variant v3: Cengage Theory p. 30]} In an amalgam concentration cell, why is $E^\circ_{\text{cell}} = 0.00\,\text{V}$ identically?
\end{itemize}
\end{problembox}

\section{Subtopic 4.2: Electrolyte Concentration Cells Without Transference}

\begin{problembox}
\textbf{Q.4.2} \hfill \textbf{[Primary Source: GRB Practice Problem 38, p. 49 / PRSN Ex-II Q22]}
\label{q:qb-c4-02}

The EMF of the concentration cell without transference:
\begin{equation*}
    \ce{Ag(s)} \mid \ce{AgNO3}(0.01\,\text{M}) \parallel \ce{AgNO3}(0.1\,\text{M}) \mid \ce{Ag(s)}
\end{equation*}
is $0.0575\,\text{V}$ at $25^\circ\text{C}$.
\begin{enumerate}[label=(\alph*)]
    \item Calculate the theoretical EMF assuming complete ionization and unit activity coefficients.
    \item If the mean activity coefficient of $0.01\,\text{M } \ce{AgNO3}$ is $\gamma_1 = 0.89$, determine the mean activity coefficient $\gamma_2$ of the $0.1\,\text{M}$ solution.
\end{enumerate}

\tcbline
\textbf{\color{accentamber}Cross-Book Redundant / Identical Variants:}
\begin{itemize}[leftmargin=*,itemsep=2pt]
    \item \textbf{[Variant v1: Neeraj Kumar Part 1 Ex-I Q98, p. 9]} The EMF of the cell $\ce{Zn} \mid \ce{Zn^2+}(0.001\,\text{M}) \parallel \ce{Zn^2+}(0.1\,\text{M}) \mid \ce{Zn}$ at $298\,\text{K}$ is: (a) $0.0591\,\text{V}$ (b) $0.0296\,\text{V}$ (c) $0.1182\,\text{V}$ (d) zero.
    \item \textbf{[Variant v2: Narendra Avasthi Level-1 Q135, p. 19]} If the concentration of both half-cells in an electrolyte concentration cell is doubled, the cell EMF will: (a) double (b) halve (c) remain unchanged (d) become zero.
    \item \textbf{[Variant v3: Cengage DPP 3.3 Q18, p. 6]} For a concentration cell reversible with respect to anions, $\ce{Ag, AgCl(s)} \mid \ce{KCl}(a_1) \parallel \ce{KCl}(a_2) \mid \ce{AgCl(s), Ag}$, the cell EMF is positive when: (a) $a_1 > a_2$ (b) $a_2 > a_1$ (c) $a_1 = a_2$ (d) always negative.
\end{itemize}
\end{problembox}

\section{Subtopic 4.3: Concentration Cells With Transference \& Liquid Junction Potential ($E_{lj}$)}

\begin{problembox}
\textbf{Q.4.3} \hfill \textbf{[Primary Source: Atkins Topic 6C.2 / Neeraj Kumar Part 1 Ex-II Sec C Passage 1]}
\label{q:qb-c4-03}

For the cell with direct liquid junction:
\begin{equation*}
    \ce{Pt, H2}(1\,\text{atm}) \mid \ce{HCl}(a_1 = 0.01) : \ce{HCl}(a_2 = 0.1) \mid \ce{H2}(1\,\text{atm})\ce{, Pt}
\end{equation*}
The EMF with transference is measured to be $E_{wt} = 0.0955\,\text{V}$ at $25^\circ\text{C}$.
\begin{enumerate}[label=(\alph*)]
    \item Derive the expression $E_{wt} = 2t_- \frac{RT}{F}\ln \frac{a_2}{a_1}$.
    \item Calculate the transport number of the chloride ion $t_{\ce{Cl-}}$ and the proton $t_{\ce{H+}}$.
    \item Calculate the liquid junction potential $E_{lj}$ at the interface.
\end{enumerate}

\tcbline
\textbf{\color{accentamber}Cross-Book Redundant / Identical Variants:}
\begin{itemize}[leftmargin=*,itemsep=2pt]
    \item \textbf{[Variant v1: Pearson Ex-II Sec A Q52, p. 20]} The liquid junction potential for a $1:1$ electrolyte boundary between activities $a_1$ and $a_2$ is given by: (a) $(t_+ - t_-)\frac{RT}{F}\ln(a_2/a_1)$ (b) $(2t_- - 1)\frac{RT}{F}\ln(a_2/a_1)$ (c) $(1 - 2t_+)\frac{RT}{F}\ln(a_2/a_1)$ (d) All of these are equivalent.
    \item \textbf{[Variant v2: GRB  12.26 Solved Ex 42, p. 33]} Explain why the liquid junction potential is virtually eliminated when a saturated \ce{KCl} bridge is inserted between two solutions.
    \item \textbf{[Variant v3: Narendra Avasthi Level-3 Passage 2, p. 37]} In which case is the liquid junction potential negative for $a_2 > a_1$? (a) $t_+ > t_-$ (b) $t_- > t_+$ (c) $t_+ = t_-$ (d) always positive.
\end{itemize}
\end{problembox}
